\documentclass{article} 
\usepackage[T1]{fontenc} 
%\usepackage[utf8]{inputenc} 
\usepackage{graphicx}
\usepackage{hyperref} 
\usepackage{float}
%\usepackage{multirow} 
%\usepackage{tabularx} 
\usepackage{color} 
%\usepackage{textcomp} 
%\usepackage{tipa}
%\usepackage{amsmath} 
%\usepackage{amssymb} 
%\usepackage{amsfonts}
%\usepackage{amsxtra} 
%\usepackage{wasysym} 
%\usepackage{isomath} 
%\usepackage{mathtools} 
%\usepackage{txfonts} 
%\usepackage{upgreek} 
%\usepackage{enumerate} 
%\usepackage{tensor} 
%\usepackage{pifont} 
%\usepackage{ulem} 
%\usepackage{xfrac} 
%\usepackage{arydshln} 
%\def\spanishoptions{mexico}
%\usepackage[english,spanish]{babel}
%\definecolor{color-DDDDDD}{rgb}{0.87,0.87,0.87}
%\definecolor{color-F9F9F9}{rgb}{0.98,0.98,0.98}
%\definecolor{white}{rgb}{1,1,1}
\definecolor{yellow}{rgb}{1,1,0}
%\definecolor{color-0F4761}{rgb}{0.06,0.28,0.38}
%los dos de abajo para la tabla de estado del arte
\usepackage{tabularx}
\usepackage{array}
%\usepackage[top=1.9cm, bottom=3cm, left=1.3cm, right=1.3cm]{geometry}
\usepackage[margin=2cm]{geometry}

\begin{document}


%Es un termino ambiguo relacionado al fine tuning o prompt engineering: especificar conforme se avance en el desarrollo Es un termino ambiguo relacionado al fine tuning o prompt engineering: especificar conforme se avance en el desarrollo

%\tableofcontents

%Análisis estadístico sobre números de teléfono en México, estadísticas sobre fraude que demuestren la veracidad de la problemática   Análisis estadístico sobre números de teléfono en México, estadísticas sobre fraude que demuestren la veracidad de la problemática   Análisis estadístico sobre números de teléfono en México, estadísticas sobre fraude que demuestren la veracidad de la problemática        Análisis estadístico sobre números de teléfono en México, estadísticas sobre fraude que demuestren la veracidad de la problemática

% Se puede hablar también sobre lo necesario que es el uso de internet y redes sociales hoy dia Se puede hablar también sobre lo necesario que es el uso de internet y redes sociales hoy dia

%Acceso a internet, uso de redes sociales, redes sociales laborales Acceso a internet, uso de redes sociales, redes sociales laborales

%Aquí se puede hablar conforme al por que es local el modelo considerando la cobertura que en ocasiones no es formidable y en la misma ciudad en edificios aveces se queda desconectado totalmente sin estar en zonas rurales Aquí se puede hablar conforme al por que es local el modelo considerando la cobertura que en ocasiones no es formidable y en la misma ciudad en edificios aveces se queda desconectado totalmente sin estar en zonas rurales

\section{\label{ref-001}Introducción}

El crecimiento exponencial en el uso de tecnologías móviles y plataformas de comunicación digital ha transformado la manera en que interactúan las personas. Según datos de la ENDUTIH 2023, en México el 81.2\,\% de la población utiliza teléfono celular y el 96.1\,\% de los usuarios de internet en la Ciudad de México lo hace para comunicarse, principalmente a través de redes sociales y aplicaciones de mensajería. Esta adopción masiva ha traído consigo grandes beneficios, pero también ha abierto nuevas oportunidades para la comisión de delitos como el \textit{catfishing} y la extorsión digital.

La facilidad con la que se pueden crear perfiles falsos y establecer contacto anónimo en entornos digitales ha hecho que estos delitos proliferen, afectando especialmente a los jóvenes, quienes presentan la mayor penetración de internet (96.7\,\% en el grupo de 18 a 24 años). Además, cifras del INEGI (ENVIPE 2023) reflejan un repunte en la tasa de incidencia de fraudes y extorsiones, situando a este fenómeno como una amenaza persistente en el ámbito digital.

\subsubsection{Personas usuarias de tecnologías de la información y comunicación}

\begin{figure}[H]
    \centering
    \includegraphics[width=1\linewidth]{imagenes/grafica1.png}
    \caption{Porcentaje de usuarios de internet, telefonía móvil y computadora en México, 2023. Adaptado de \textit{Encuesta Nacional sobre Disponibilidad y Uso de Tecnolog{\'{i}}as de la Informaci{\'{o}}n en los Hogares (ENDUTIH) 2023}, por Instituto Nacional de Estad{\'{i}}stica y Geograf{\'{i}}a, 2024, \url{https://www.inegi.org.mx/programas/endutih/2023/\#informacion_general}.
    {\textcopyright} INEGI}
    \label{fig:1}
\end{figure}

\begin{figure}[H]
    \centering
    \includegraphics[width=1\linewidth]{imagenes/grafica2.png}
    \caption{Personas usuarias de internet seg{{\'{u}}}n grupos de edad del 2020 al 2023\newline
    Acceso a internet por entidad federativa, 2023. Tomado de \textit{Encuesta Nacional sobre Disponibilidad y Uso de Tecnolog{\'{i}}as de la Informaci{\'{o}}n en los Hogares (ENDUTIH) 2023: Comunicado de prensa n{{\'{u}}}m. 161/24}, por Instituto Nacional de Estad{\'{i}}stica y Geograf{\'{i}}a, 2024, \url{https://www.inegi.org.mx/contenidos/saladeprensa/boletines/2024/ENDUTIH/ENDUTIH_23.pdf}. {\textcopyright} INEGI}
    \label{fig:2}
\end{figure}

Datos de la ENDUTIH 2024 muestran que los jóvenes de \textbf{18 a 24 años} presentan la mayor penetración de Internet en México (96.7 \%), lo cual coincide con su alta participaci{\'{o}}n en redes sociales y aplicaciones de citas. Esta combinaci{\'{o}}n los convierte en un grupo con mayor exposici{\'{o}}n a \textbf{estafas rom{\'{a}}nticas y catfishing}

\subsubsection{Tasa de incidencia delictiva por cada cien mil habitantes por tipo de delito.}

% Tabla de incidencia delictiva
\begin{table}[H]
  \centering
  \caption{Tasa de incidencia delictiva por cada cien mil habitantes por tipo de delito.}
  \label{tab:incidencia}
  \begin{tabular}{@{}lcrrrr@{}}
    \toprule
    Año & Tipo de delito & Estimaciones puntuales & Coef. de variación (\%) & Límite inferior & Límite superior \\
    \midrule
    2023 & Fraude & 6.962 & 2.4 & 6.687 & 7.238 \\
         & Extorsión & 5.213 & 3.2 & 4.940 & 5.486 \\
    2022 & Fraude & 5.770 & 2.5 & 5.530 & 6.009 \\
         & Extorsión & 5.056 & 3.0 & 4.803 & 5.308 \\
    2020 & Fraude & 5.904 & 2.5 & 5.660 & 6.148 \\
         & Extorsión & 5.160 & 3.0 & 4.909 & 5.412 \\
    2018 & Fraude & 5.397 & 3.0 & 5.133 & 5.661 \\
         & Extorsión & 6.542 & 2.7 & 6.248 & 6.836 \\
    2016 & Fraude & 4.656 & 2.9 & 4.433 & 4.879 \\
         & Extorsión & 8.945 & 2.3 & 8.606 & 9.283 \\
    2014 & Fraude & 4.255 & 4.6 & 3.933 & 4.576 \\
         & Extorsión & 9.850 & 2.8 & 9.404 & 10.295 \\
    2012 & Fraude & 3.319 & 4.0 & 3.101 & 3.536 \\
         & Extorsión & 7.585 & 3.0 & 7.216 & 7.954 \\
    2010 & Fraude & 2.658 & 4.3 & 2.469 & 2.848 \\
         & Extorsión & 7.239 & 4.3 & 6.727 & 7.750 \\
    \bottomrule
  \end{tabular}
  \footnotesize{\\ \centering Tasa de incidencia delictiva por tipo de delito en México, 2023. Tomado de \textit{Encuesta Nacional de Victimizaci{\'{o}}n y Percepción sobre Seguridad P{{\'{u}}}blica (ENVIPE) \textendash{} Incidencia delictiva por tipo de delito}, por Instituto Nacional de Estad{\'{i}}stica y Geograf{\'{i}}a, 2024, \url{https://www.inegi.org.mx/app/tabulados/interactivos/?px=ENVIPE_INCIDENCIA_TIPO_DELITO&bd=ENVIPE}. {\textcopyright} INEGI }
\end{table}

\subsubsection{Servicios con los que cuenta la poblaci{\'{o}}n}

Seg{{\'{u}}}n la encuesta ``Hogares con equipamiento de tecnolog{\'{i}}a de informaci{\'{o}}n y comunicaciones, por entidad federativa, seg{{\'{u}}}n tipo de tecnolog{\'{i}}a, 2023'' en los hogares de la ciudad de M{\'{e}}xico se cuentan con los servicios de Internet 89.5\%, telefon{\'{i}}a 98.9\%

Instituto Nacional de Estad{\'{i}}stica y Geograf{\'{i}}a. (2024). \textit{Encuesta Nacional sobre Disponibilidad y Uso de Tecnolog{\'{i}}as de la Informaci{\'{o}}n en los Hogares (ENDUTIH) 2023 \textendash{} Tabulados de resultados} {[}Archivo de Excel{]}. \url{https://www.inegi.org.mx/contenidos/programas/endutih/2023/tabulados/2023_h2ed110.xlsx}

\subsubsection{Prop{\'{o}}sitos del uso de internet}

Seg{{\'{u}}}n la base de datos ``Usuarios de internet, por entidad federativa, seg{{\'{u}}}n principales usos, 2023'', El 96,1 \% de los 7,897,062 de personas encuestadas en la ciudad de M{\'{e}}xico , utilizan el {\textquestiondown} para comunicarse, el 90,7 \% para redes sociales, el 37.7 \% para realizar operaciones bancarias, el 1.3 para pagos con tarjetas de prepago

Instituto Nacional de Estad{\'{i}}stica y Geograf{\'{i}}a. (2023). \textit{ENDUTIH 2023: Tabulados {[}Archivo de Excel{]}}. \url{https://www.inegi.org.mx/contenidos/programas/endutih/2023/tabulados/2023_u2ed244.xlsx}

\section{\label{ref-002}Planteamiento del problema%Edgar Diego Pina VargasEP2025-05-08T20:09:00Z0 Edgar Diego Pina VargasEP Edgar Diego Pina VargasEP Edgar Diego Pina Vargas Edgar Diego Pina Vargas EP EP 2025-05-08T20:09:00Z 2025-05-08T20:09:00Z 0 0 0   Como este ahorita el planteamiento solo aborda el problema pero no como se piensa solucionar. En el planteamiento del problema se aborda primeramente el problema y posteriormente se puede hablar sobre los posibles enfoques con los que se podría solucionar el problema y dentro de estos enfoques   Hablar sobre las consecuencias económicas, psicológicas, de privacidad   Como este ahorita el planteamiento solo aborda el problema pero no como se piensa solucionar. En el planteamiento del problema se aborda primeramente el problema y posteriormente se puede hablar sobre los posibles enfoques con los que se podría solucionar el problema y dentro de estos enfoques    Como este ahorita el planteamiento solo aborda el problema pero no como se piensa solucionar. En el planteamiento del problema se aborda primeramente el problema y posteriormente se puede hablar sobre los posibles enfoques con los que se podría solucionar el problema y dentro de estos enfoques         Como este ahorita el planteamiento solo aborda el problema pero no como se piensa solucionar. En el planteamiento del problema se aborda primeramente el problema y posteriormente se puede hablar sobre los posibles enfoques con los que se podría solucionar el problema y dentro de estos enfoques    Hablar sobre las consecuencias económicas, psicológicas, de privacidad   Hablar sobre las consecuencias económicas, psicológicas, de privacidad        Hablar sobre las consecuencias económicas, psicológicas, de privacidad
}

Hoy en d{\'{i}}a el crecimiento de internet y aspectos multifactoriales han dado como producto un crecimiento de las estafas en l{\'{i}}nea y diversos tipos de extorci{\'{o}}n a trav{\'{e}}s de las redes sociales, en la actualidad este tipo de fen{\'{o}}menos son principalmente efectuados dentro de los servicios de mensajer{\'{i}}a instant{\'{a}}nea, los cuales se convierten en una oportunidad que tienen los ciberdelincuentes para establecer contacto con diferentes usuarios a lo largo del mundo de manera an{\'{o}}nima, la falta de evidencia documentada sobre las estafas y casos de extorci{\'{o}}n en medios digitales tambi{\'{e}}n forma parte principal de los problemas en la agilizaci{\'{o}}n de los procesos que enfrentan las autoridades judiciales.

\section{Hip{\'{o}}tesis %Edgar Diego Pina VargasEP2025-05-08T20:05:00Z0 Edgar Diego Pina VargasEP Edgar Diego Pina VargasEP Edgar Diego Pina Vargas Edgar Diego Pina Vargas EP EP 2025-05-08T20:05:00Z 2025-05-08T20:05:00Z 0 0 0   Redactar una hipótesis hablar sobre el enfoque de la solución con la hipótesis de que el usuario no caiga en el fraude, 2da hipótesis la parte de reportes, la hipótesis por tanto permite hablar del enfoque con el que se atacara el problema   Redactar una hipótesis hablar sobre el enfoque de la solución con la hipótesis de que el usuario no caiga en el fraude, 2da hipótesis la parte de reportes, la hipótesis por tanto permite hablar del enfoque con el que se atacara el problema   Redactar una hipótesis hablar sobre el enfoque de la solución con la hipótesis de que el usuario no caiga en el fraude, 2da hipótesis la parte de reportes, la hipótesis por tanto permite hablar del enfoque con el que se atacara el problema        Redactar una hipótesis hablar sobre el enfoque de la solución con la hipótesis de  que el usuario no caiga en el fraude, 2da hipótesis la parte de reportes,  que el usuario no caiga en el fraude, 2da hipótesis la parte de reportes,  la hipótesis por tanto permite hablar del enfoque con el que se atacara el problema
%Edgar Diego Pina VargasEP2025-05-09T21:50:00Z0 Edgar Diego Pina VargasEP Edgar Diego Pina VargasEP Edgar Diego Pina Vargas Edgar Diego Pina Vargas EP EP 2025-05-09T21:50:00Z 2025-05-09T21:50:00Z 0 0 0 Desarrollar la hipótesis que se tiene sobre “un sistema que alerta sobre mensajes con contenido de catphising prevendría el que el usuario sea victima del fraude” Desarrollar la hipótesis que se tiene sobre “un sistema que alerta sobre mensajes con contenido de catphising prevendría el que el usuario sea victima del fraude” Desarrollar la hipótesis que se tiene sobre “un sistema que alerta sobre mensajes con contenido de catphising prevendría el que el usuario sea victima del fraude” La segunda hipótesis sobre “el reporteador de evidencias y su apoyo en el proceso legal“, cuidar el afirmar como evidencias totales, se podría hablar de parciales y fundamentar desde el alcance tecnológico con el nivel de pertenencia de datos y aspectos legales sobre lo invasivo de los datos La segunda hipótesis sobre “el reporteador de evidencias y su apoyo en el proceso legal“, cuidar el afirmar como evidencias totales, se podría hablar de parciales y fundamentar desde el alcance tecnológico con el nivel de pertenencia de datos y aspectos legales sobre lo invasivo de los datos La segunda hipótesis sobre “el reporteador de evidencias y su apoyo en el proceso legal“, cuidar el afirmar como evidencias totales, se podría hablar de parciales y fundamentar desde el alcance tecnológico con el nivel de pertenencia de datos y aspectos legales sobre lo invasivo de los datos
}

\section{\label{ref-003}Objetivo general \ding{52}}

\colorbox{yellow}{{\textquestiondown}Dise{\~{n}}ar/implementar?} un sistema de filtrado de mensajes de texto fraudulentos dentro de la clasificaci{\'{o}}n de extorci{\'{o}}n y catphishing mediante un LLM refinado mediante t{\'{e}}cnicas de prompt engineering y fine tuning%Edgar Diego Pina VargasEP2025-04-30T11:15:00Z0 Edgar Diego Pina VargasEP Edgar Diego Pina VargasEP Edgar Diego Pina Vargas Edgar Diego Pina Vargas EP EP 2025-04-30T11:15:00Z 2025-04-30T11:15:00Z 0 0 0 Fine tuning a nivel de prompt engineering no como técnica de entrenamiento Fine tuning a nivel de prompt engineering no como técnica de entrenamiento


\section{\label{ref-004}Objetivos espec{\'{i}}ficos%Edgar Diego Pina VargasEP2025-05-09T21:26:00Z0 Edgar Diego Pina VargasEP Edgar Diego Pina VargasEP Edgar Diego Pina Vargas Edgar Diego Pina Vargas EP EP 2025-05-09T21:26:00Z 2025-05-09T21:26:00Z 0 0 0 Simplificar la redacción de los objetivos priorizando los action verbs, ya en el marco teorico se puede hablar mas detenidamente sobre las características profundoas de los modelos y todo lo contemplado en profundidad tecnica Simplificar la redacción de los objetivos priorizando los action verbs, ya en el marco teorico se puede hablar mas detenidamente sobre las características profundoas de los modelos y todo lo contemplado en profundidad tecnica
}

Generar datos sint{\'{e}}ticas puntualmente ejemplos de catphising y extorci{\'{o}}n mediante uso de modelos de lenguaje preentrenados para realizar pruebas de desempe{\~{n}}o a los diferentes modelos de lenguaje contemplados y definir el modelo base a utilizar en la solucion

  	Investigaci{\'{o}}n an{\'{a}}lisis y selecci{\'{o}}n de {\ldots}

Definir el/los prompts que entreguen un mejor resultado en el modelo seleccionado mediante pruebas considerando la efectividad al detectar correctamente el catphising y/o extorci{\'{o}}n en los ejemplos suministrados para registrarlos y posteriormente implementarlos en la aplicaci{\'{o}}n de mensajer{\'{i}}a%Edgar Diego Pina VargasEP2025-05-09T21:46:00Z0 Edgar Diego Pina VargasEP Edgar Diego Pina VargasEP Edgar Diego Pina Vargas Edgar Diego Pina Vargas EP EP 2025-05-09T21:46:00Z 2025-05-09T21:46:00Z 0 0 0 Como el tecnisismo prompt puede tener colision, investigar como poner notas al pie de pagina para distinguir a que se refiere en nuestro ámbito o directamente aconsejar visitar la sección de glosario de términos, analizar esto no solo para prompt sino para tantos se requieran (conviene mejor una redaccion glocal que recomiende visitar el glosario) Como el tecnisismo prompt puede tener colision, investigar como poner notas al pie de pagina para distinguir a que se refiere en nuestro ámbito o directamente aconsejar visitar la sección de glosario de términos, analizar esto no solo para prompt sino para tantos se requieran (conviene mejor una redaccion glocal que recomiende visitar el glosario)


Seguir uno de los dos enfoques para llevar a cabo el fine tuning en medida de haber conseguido resultados de identificaci{\'{o}}n solidos con el modelo seleccionado y el conjunto de prompts refinados

Implementar una aplicaci{\'{o}}n de mensajer{\'{i}}a nativa en Andoid %Edgar Diego Pina VargasEP2025-05-09T22:24:00Z0 Edgar Diego Pina VargasEP Edgar Diego Pina VargasEP Edgar Diego Pina Vargas Edgar Diego Pina Vargas EP EP 2025-05-09T22:24:00Z 2025-05-09T22:24:00Z 0 0 0 La selección de Android debe ser justificada también, puede ser en marco teorico o en el desarrollo, en dhica parte habría que resaltar las ventajas y desventajas de cada tecnología para justificar la elección de la misma, se puede considerar la experiencia de los autores con la tecnología pero es mejor desde las características técnicas y de servicio como los que incluyen soporte técnico en contraste con el open sourse La selección de Android debe ser justificada también, puede ser en marco teorico o en el desarrollo, en dhica parte habría que resaltar las ventajas y desventajas de cada tecnología para justificar la elección de la misma, se puede considerar la experiencia de los autores con la tecnología pero es mejor desde las características técnicas y de servicio como los que incluyen soporte técnico en contraste con el open sourse un aspecto puede ser factor socioeconómico de México un aspecto puede ser factor socioeconómico de México
mediante Android Studio para conseguir tener un primer prototipo funcional

Migrar el modelo refinado mediante fine tuning al ecosistema de Andoid mediante la tecnolog{\'{i}}a LiteRT para ejecutarlo como un servicio local en el dispositivo

\section{\label{ref-005}Justificaci{\'{o}}n%Edgar Diego Pina VargasEP2025-05-09T22:42:00Z0 Edgar Diego Pina VargasEP Edgar Diego Pina VargasEP Edgar Diego Pina Vargas Edgar Diego Pina Vargas EP EP 2025-05-09T22:42:00Z 2025-05-09T22:42:00Z 0 0 0 La justificación es puramente el por que se tiene que hacer, la parte de justificar las tecnologías e infraestructura esta mas de la mano del marco teorico o el desarrollo, la parte de la falta de evidencias también puede gihurar en la justificacion La justificación es puramente el por que se tiene que hacer, la parte de justificar las tecnologías e infraestructura esta mas de la mano del marco teorico o el desarrollo, la parte de la falta de evidencias también puede gihurar en la justificacion
	%Edgar Diego Pina VargasEP2025-05-09T22:35:00Z0 Edgar Diego Pina VargasEP Edgar Diego Pina VargasEP Edgar Diego Pina Vargas Edgar Diego Pina Vargas EP EP 2025-05-09T22:35:00Z 2025-05-09T22:35:00Z 0 0 0 Altamente Relacionado con la hipótesis pues en la justificación se hablara de lo que beneficiaria el desarrollo de la solución planteada desde la hipotesis Altamente Relacionado con la hipótesis pues en la justificación se hablara de lo que beneficiaria el desarrollo de la solución planteada desde la hipotesis
}

\section{\label{ref-006}Marco Te{\'{o}}rico%Edgar Diego Pina VargasEP2025-05-08T20:01:00Z0 Edgar Diego Pina VargasEP Edgar Diego Pina VargasEP Edgar Diego Pina Vargas Edgar Diego Pina Vargas EP EP 2025-05-08T20:01:00Z 2025-05-08T20:01:00Z 0 0 0   En marco teórico entrar en aspectos técnicos más profundos por sobre la redacción de los objetivos específicos   En marco teórico entrar en aspectos técnicos más profundos por sobre la redacción de los objetivos específicos   En marco teórico entrar en aspectos técnicos más profundos por sobre la redacción de los objetivos específicos        En marco teórico entrar en aspectos técnicos más profundos por sobre la redacción de los objetivos específicos   Incluir definición de lenguaje en una sección del marco teórico, identificándolo como un ente vivo y cambiante a lo largo del tiempo, que posee regionalismos, jergas posteriormente relacionar estas características del lenguaje con las características que debe del tener el llm base para el filtro se argumenta la complejidad de la existencia de un modelo con tales modismos y posteriormente se justifica la necesidad de refinar con fine tuning a nivel de prompts  Aquí añade los aspectos técnicos que habrá que cubrir, considera como guía la sección de la propuesta de solución que habla mas en profundidad sobre los aspectos técnicos   Incluir definición de lenguaje en una sección del marco teórico, identificándolo como un ente vivo y cambiante a lo largo del tiempo, que posee regionalismos, jergas posteriormente relacionar estas características del lenguaje con las características que debe del tener el llm base para el filtro se argumenta la complejidad de la existencia de un modelo con tales modismos y posteriormente se justifica la necesidad de refinar con fine tuning a nivel de prompts   Incluir definición de lenguaje en una sección del marco teórico, identificándolo como un ente vivo y cambiante a lo largo del tiempo, que posee regionalismos, jergas posteriormente relacionar estas características del lenguaje con las características que debe del tener el llm base para el filtro se argumenta la complejidad de la existencia de un modelo con tales modismos y posteriormente se justifica la necesidad de refinar con fine tuning a nivel de prompts        Incluir definición de lenguaje en una sección del marco teórico, identificándolo como un ente vivo y cambiante a lo largo del tiempo, que posee regionalismos, jergas posteriormente relacionar estas características del lenguaje con las características que debe del tener el llm base para el filtro se argumenta la complejidad de la existencia de un modelo con tales modismos y posteriormente se justifica la necesidad de refinar con fine tuning a nivel de prompts      Aquí añade los aspectos técnicos que habrá que cubrir, considera como guía la sección de la propuesta de solución que habla mas en profundidad sobre los aspectos técnicos   Aquí añade los aspectos técnicos que habrá que cubrir, considera como guía la sección de la propuesta de solución que habla mas en profundidad sobre los aspectos técnicos        Aquí añade los aspectos técnicos que habrá que cubrir, considera como guía la sección de la propuesta de solución que habla mas en profundidad sobre los aspectos técnicos
%Edgar Diego Pina VargasEP2025-05-09T21:14:00Z0 Edgar Diego Pina VargasEP Edgar Diego Pina VargasEP Edgar Diego Pina Vargas Edgar Diego Pina Vargas EP EP 2025-05-09T21:14:00Z 2025-05-09T21:14:00Z 0 0 0 Línea 2:  Línea 2:  DEP 2 (Entrega: 09-may-2025, 17:30 horas) DEP 2 (Entrega: 09-may-2025, 17:30 horas) (20\%) MARCO TEÓRICO (20\%) MARCO TEÓRICO - (10\%) Citas - (10\%) Citas - (10\%) Conceptos, definiciones y teorías que requiere el lector para entender su - (10\%) Conceptos, definiciones y teorías que requiere el lector para entender su trabajo. trabajo. (80\%) PROPUESTA DE DESARROLLO (80\%) PROPUESTA DE DESARROLLO - (30\%) Diagrama a bloques o algoritmo(s) de la propuesta de desarrollo - (30\%) Diagrama a bloques o algoritmo(s) de la propuesta de desarrollo - (30\%) Explicación con texto del diagrama a bloques o algoritmo(s) de la - (30\%) Explicación con texto del diagrama a bloques o algoritmo(s) de la propuesta de desarrollo propuesta de desarrollo - (10\%) Citas - (10\%) Citas - (10\%) Explicaciones adicionales sobre la propuesta de desarrollo - (10\%) Explicaciones adicionales sobre la propuesta de desarrollo
}

\textbf{\textit{Que son los modelos largos de lenguaje LLM y cual es su utilidad}}

Los modelos largos de lenguajes o modelos de lenguaje de gran tama{\~{n}}o, son modelos de aprendizaje profundo que son pre entrenados con grandes cantidades de datos 

Poseen una estructura fundamentada en los transformadores, un tipo de arquitectura introducido en el art{\'{i}}culo "Attention Is All You Need", este tipo de arquitectura se basa exclusivamente en mecanismos de atenci{\'{o}}n, eliminando la necesidad de recurrencias o convoluciones. Los transformadores son un conjunto de redes neuronales que cuenta con un codificador y decodificador con capacidades de autocorrecci{\'{o}}n estos extraen significados de una secuencia de texto y comprenden las relaciones entre las palabras y las frases que contiene. {[}3{]}{[}1{]}

Los LLMS pueden ser empleados para llevar a cabo tareas de diversas indoles como pueden ser, responder preguntas, resumir documentos, traducir idiomas, completar oraciones, automatizar la creaci{\'{o}}n de contenido y cambiar la forma en que las personas utilizan los motores de b{{\'{u}}}squeda y los asistentes virtuales.

Particularmente en el presente trabajo nos centraremos particularmente en la tarea de clasificaci{\'{o}}n de textos

\textbf{\textit{LLMs Como clasificadores de textos}}
\begin{itemize}
\item[►] \foreignlanguage{spanish}{Los modelos de lenguaje de gran escala (LLMs, por sus siglas en ingl{\'{e}}s) han transformado la forma en que se abordan las tareas de procesamiento de lenguaje natural (PLN), especialmente en el {\'{a}}mbito de la clasificaci{\'{o}}n de textos. Tradicionalmente, esta tarea depend{\'{i}}a de modelos supervisados entrenados con grandes vol{{\'{u}}}menes de datos etiquetados, los cuales requer{\'{i}}an etapas expl{\'{i}}citas de extracci{\'{o}}n de caracter{\'{i}}sticas y entrenamiento. En contraste, }\foreignlanguage{spanish}{\colorbox{yellow}{los LLMs permiten realizar tareas de clasificaci{\'{o}}n incluso en entornos con pocos o ning{{\'{u}}}n dato etiquetado, a trav{\'{e}}s de t{\'{e}}cnicas como el }}\foreignlanguage{spanish}{\colorbox{yellow}{\textit{zero-shot learning}}}\foreignlanguage{spanish}{\colorbox{yellow}{ y el }}\foreignlanguage{spanish}{\colorbox{yellow}{\textit{few-shot learning}}}\foreignlanguage{spanish}{\colorbox{yellow}{.}}
\item[►] \foreignlanguage{spanish}{En el art{\'{i}}culo ``}\foreignlanguage{spanish}{\colorbox{yellow}{Large Language Models as Zero-Shot Classifiers: An Empirical Study}}\foreignlanguage{spanish}{'' (arXiv:2312.01044), se realiz{\'{o}} una evaluaci{\'{o}}n emp{\'{i}}rica del desempe{\~{n}}o de distintos LLMs como clasificadores sin entrenamiento adicional, usando {{\'{u}}}nicamente instrucciones (prompts) dise{\~{n}}adas cuidadosamente. Entre los hallazgos m{\'{a}}s relevantes destaca que, }\foreignlanguage{spanish}{\colorbox{yellow}{\textbf{\textit{aunque estos modelos no superan sistem{\'{a}}ticamente a clasificadores supervisados entrenados sobre datos espec{\'{i}}ficos, ofrecen una ventaja considerable en contextos donde no se dispone de datos etiquetados o los recursos computacionales para entrenamiento son limitados. }}}}\foreignlanguage{spanish}{Asimismo, se demostr{\'{o}} que el dise{\~{n}}o del prompt es determinante: un cambio sutil en la redacci{\'{o}}n de la instrucci{\'{o}}n puede afectar significativamente la precisi{\'{o}}n del modelo.}
\end{itemize}
\foreignlanguage{spanish}{Para el proyecto aqu{\'{i}} concebido, que busca detectar autom{\'{a}}ticamente patrones de catfishing y extorsi{\'{o}}n en mensajes dentro de aplicaciones de mensajer{\'{i}}a, el uso de LLMs como clasificadores resulta particularmente pertinente. Dado que estos delitos evolucionan constantemente y pueden adoptar m{{\'{u}}}ltiples formas ling{\"{u}}{\'{i}}sticas, los modelos tradicionales entrenados sobre conjuntos de datos fijos podr{\'{i}}an mostrar limitaciones en su capacidad de generalizaci{\'{o}}n. En cambio, un LLM permite adaptar din{\'{a}}micamente su capacidad de clasificaci{\'{o}}n mediante modificaciones en los prompts, o bien mediante t{\'{e}}cnicas de }\foreignlanguage{spanish}{\textit{fine tuning}}\foreignlanguage{spanish}{ si se desea aumentar la precisi{\'{o}}n en un dominio espec{\'{i}}fico.}

\foreignlanguage{spanish}{\textbf{\textit{Definici{\'{o}}n formal del lenguaje}}}
\begin{itemize}
\item[►] \foreignlanguage{spanish}{Desde la perspectiva de la ling{\"{u}}{\'{i}}stica formal, el lenguaje se concibe como un sistema estructurado de signos que permite la comunicaci{\'{o}}n entre individuos. Este sistema se caracteriza por una gram{\'{a}}tica que define las reglas para la formaci{\'{o}}n de expresiones v{\'{a}}lidas dentro del idioma. }
\item[►] \foreignlanguage{spanish}{La ling{\"{u}}{\'{i}}stica formal se enfoca en describir y analizar las estructuras gramaticales del lenguaje, utilizando herramientas matem{\'{a}}ticas y l{\'{o}}gicas para modelar fen{\'{o}}menos ling{\"{u}}{\'{i}}sticos. Este enfoque permite una comprensi{\'{o}}n precisa de c{\'{o}}mo se construyen y entienden las oraciones en un idioma, facilitando el desarrollo de modelos computacionales para el procesamiento del lenguaje natural.}
\end{itemize}
\foreignlanguage{spanish}{\textbf{\textit{La variabilidad en el espa{\~{n}}ol y las jergas empleadas en la extorsi{\'{o}}n y el catfishing}}}
\begin{itemize}
\item[►] \foreignlanguage{spanish}{El espa{\~{n}}ol, como lengua hablada en m{{\'{u}}}ltiples pa{\'{i}}ses y regiones, presenta una notable variabilidad dialectal y l{\'{e}}xica. Esta diversidad se manifiesta en diferencias de vocabulario, pronunciaci{\'{o}}n y uso de expresiones idiom{\'{a}}ticas, influenciadas por factores geogr{\'{a}}ficos, hist{\'{o}}ricos y sociales.}
\item[►] \foreignlanguage{spanish}{Un estudio reciente analiz{\'{o}} mensajes p{{\'{u}}}blicos de Twitter geolocalizados en 26 pa{\'{i}}ses de habla hispana, identificando variaciones significativas en el uso del lenguaje. Estas diferencias incluyen el empleo de t{\'{e}}rminos y expresiones particulares de cada regi{\'{o}}n, as{\'{i}} como la adopci{\'{o}}n de jergas espec{\'{i}}ficas en contextos sociales determinados.}
\item[►] \foreignlanguage{spanish}{En el {\'{a}}mbito de la extorsi{\'{o}}n y el catfishing, los delincuentes suelen utilizar jergas y modismos locales para ganarse la confianza de sus v{\'{i}}ctimas y evitar la detecci{\'{o}}n por sistemas automatizados. Esta adaptaci{\'{o}}n ling{\"{u}}{\'{i}}stica representa un desaf{\'{i}}o para los modelos de procesamiento de lenguaje natural, que deben ser capaces de reconocer y entender estas variaciones para identificar comportamientos sospechosos de manera efectiva.}
\item[►] \foreignlanguage{spanish}{Por lo tanto, es fundamental que los sistemas de detecci{\'{o}}n de actividades fraudulentas en redes sociales incorporen mecanismos para manejar la variabilidad ling{\"{u}}{\'{i}}stica del espa{\~{n}}ol. Esto incluye el entrenamiento de modelos con datos representativos de diferentes dialectos y jergas, as{\'{i}} como la actualizaci{\'{o}}n constante de los recursos ling{\"{u}}{\'{i}}sticos utilizados para reflejar los cambios y tendencias en el uso del lenguaje.}
\end{itemize}
Large Lenguaje Models

Desarrollar este objetivo

Seleccionar un modelo de lenguaje base adecuado balanceando tama{\~{n}}o y habilidad para identificaci{\'{o}}n de catphising y extorci{\'{o}}n mediante la ejecuci{\'{o}}n de pruebas en el entorno de jypiter lab %Edgar Diego Pina VargasEP2025-05-09T21:28:00Z0 Edgar Diego Pina VargasEP Edgar Diego Pina VargasEP Edgar Diego Pina Vargas Edgar Diego Pina Vargas EP EP 2025-05-09T21:28:00Z 2025-05-09T21:28:00Z 0 0 0 Sintetizado : Investigación análisis y selección de un modelo adecuado para identificación de catphishing  Sintetizado : Investigación análisis y selección de un modelo adecuado para identificación de catphishing 


Retomar a que se refiere lo de ``modelo base'' en la parte de iobjetivos

\section{M{\'{e}}tricas de desempe{\~{n}}o de clasificadores}

\subsection{Presicion}

La precisi{\'{o}}n es la habilidad que posee un clasificador para evitar clasificar como positivo un caso negativo

\subsection{Recall}

Recall es la habilidad que tiene un clasificador para encontrar todos los casos positivos

\url{https://scikit-learn.org/stable/modules/generated/sklearn.metrics.precision_score.html\#sklearn.metrics.precision_score}

Buenas practicas para la generaci{\'{o}}n de datos sint{\'{e}}ticos (Pendiente de parafraseo y resumen)

Esta informaci{\'{o}}n generada artificialmente puede complementar o incluso reemplazar los datos del mundo real al entrenar o probar los modelos de~\href{https://www.ibm.com/mx-es/think/topics/artificial-intelligence}{inteligencia artificial}. En comparaci{\'{o}}n con los conjuntos de datos reales, que suponen desventajas en cuesti{\'{o}}n de costo, acceso, tiempo para etiquetarlos y suministro, los conjuntos de datos sint{\'{e}}ticos se pueden sintetizar mediante simulaciones por computadora o~\href{https://www.ibm.com/mx-es/think/topics/generative-model}{modelos generativos}. Por lo tanto, es m{\'{a}}s econ{\'{o}}mico producirlos a demanda en vol{{\'{u}}}menes casi ilimitados y personalizarlos en funci{\'{o}}n de las necesidades de una organizaci{\'{o}}n.

A pesar de sus beneficios, los~\href{https://research.ibm.com/blog/what-is-synthetic-data}{datos sint{\'{e}}ticos}~tambi{\'{e}}n conllevan dificultades. El proceso de generaci{\'{o}}n puede ser complejo, ya que los cient{\'{i}}ficos de datos tienen que crear datos realistas sin sacrificar su calidad y privacidad.

IBM propone estas ~ocho mejores pr{\'{a}}cticas para la generaci{\'{o}}n de datos sint{\'{e}}ticos:

1: Conocer el prop{\'{o}}sito:

Comprender por qu{\'{e}} el proyecto necesita datos sint{\'{e}}ticos y los casos de uso en los que podr{\'{i}}an ser m{\'{a}}s {{\'{u}}}tiles que los datos reales.~

Los datos artificiales pueden actuar como una forma de~\href{https://www.ibm.com/mx-es/think/topics/data-augmentation}{refuerzo de datos}. Se pueden utilizar para impulsar la diversidad de los datos, especialmente para grupos subrepresentados en el entrenamiento de~\href{https://www.ibm.com/mx-es/topics/ai-model}{modelos de IA.}~Y cuando la informaci{\'{o}}n es escasa, los datos sint{\'{e}}ticos pueden llenar los vac{\'{i}}os.

A la empresa de servicios financieros JP Morgan, por ejemplo, le result{\'{o}} dif{\'{i}}cil entrenar eficazmente modelos impulsados por IA para la detecci{\'{o}}n de fraudes debido a la falta de casos fraudulentos en comparaci{\'{o}}n con los no fraudulentos.

2: la preparaci{\'{o}}n es clave

Al preparar conjuntos de datos originales para la generaci{\'{o}}n de datos sint{\'{e}}ticos mediante algoritmos de~\href{https://www.ibm.com/mx-es/think/topics/machine-learning}{aprendizaje autom{\'{a}}tico}~(ML), aseg{{\'{u}}}rese de verificar y corregir errores, imprecisiones e incongruencias. Elimine los duplicados e ingrese los valores faltantes.

\section{Glosario de t{\'{e}}rminos%Edgar Diego Pina VargasEP2025-05-08T20:03:00Z0 Edgar Diego Pina VargasEP Edgar Diego Pina VargasEP Edgar Diego Pina Vargas Edgar Diego Pina Vargas EP EP 2025-05-08T20:03:00Z 2025-05-08T20:03:00Z 0 0 0   Presentar en orden alfabetico   Presentar en orden alfabetico   Presentar en orden alfabetico        Presentar en orden alfabetico
}

\textbf{Catphishing}: Pr{\'{a}}ctica enga{\~{n}}osa en la que una persona crea una identidad falsa en l{\'{i}}nea para establecer relaciones con otras personas, generalmente con fines de manipulaci{\'{o}}n emocional, fraude o explotaci{\'{o}}n. Es com{{\'{u}}}n en aplicaciones de citas y redes sociales, donde los estafadores se hacen pasar por alguien m{\'{a}}s para ganarse la confianza de la v{\'{i}}ctima y, en muchos casos, obtener beneficios econ{\'{o}}micos o personales.

\textbf{Ciberfraude}: Delito que implica el uso de tecnolog{\'{i}}as de la informaci{\'{o}}n y comunicaci{\'{o}}n para cometer actos fraudulentos, como el robo de identidad, estafas financieras, phishing, entre otros. El ciberfraude puede manifestarse en diversas formas, incluyendo correos electr{\'{o}}nicos enga{\~{n}}osos, sitios web falsos y suplantaci{\'{o}}n de identidad en redes sociales.

\textbf{Prompt}: Instrucci{\'{o}}n o entrada de texto proporcionada a un modelo largo de lenguaje para generar una respuesta espec{\'{i}}fica. La calidad y claridad del prompt influyen directamente en la precisi{\'{o}}n y relevancia de la salida generada por el modelo de inteligencia artificial.

\textbf{Prompt Engineering}: Disciplina que se enfoca en dise{\~{n}}ar y optimizar prompts para guiar eficazmente a modelos de lenguaje de gran escala (LLM) en la generaci{\'{o}}n de respuestas deseadas. Implica comprender c{\'{o}}mo formular instrucciones precisas que alineen las salidas del modelo con los objetivos espec{\'{i}}ficos del usuario o aplicaci{\'{o}}n.

\textbf{Zero-shot}: Es una t{\'{e}}cnica de prompt engineering que se identifica por redactar el prompt con una instrucci{\'{o}}n directa sin incluir ejemplos sobre la tarea asignada

\textbf{Few-shot}: Es una t{\'{e}}cnica de prompt engineering que es identificada por incluir ejemplos sobre la tarea que se le indica al modelo dentro del mismo prompt 

\textbf{Fine Tuning}: Proceso de adaptar un modelo de lenguaje previamente entrenado para tareas o dominios específicos mediante un entrenamiento adicional con datos relevantes. En el contexto de este proyecto, se destaca el uso de t{\'{e}}cnicas de prompt engineering para afinar el modelo sin necesidad de reentrenarlo completamente, lo que permite una personalizaci{\'{o}}n eficiente y menos costosa en t{\'{e}}rminos computacionales.

\textbf{Large language model LLM:} Son modelos de inteligencia artificial que suelen ser entrenados con grandes cantidades de informaci{\'{o}}n principalmente de textos con el fin de entender y manipular el lenguaje humano. 

\textbf{Lenguaje humano:} sistemas de signos que sirven para establecer comunicaci{\'{o}}n entre personas, especialmente mediante la palabra hablada o escrita.

Se invita a consultar la secci{\'{o}}n ``marco teorico'' para profundizar mas en cada uno de los terminos

\section{Estado del arte}

\begin{table}[H]
    \centering
    \caption{Comparación de Trabajos Relacionados}
    \begin{tabularx}{\textwidth}{|X|X|X|X|X|X|}
    %{|p{2cm}|p{2.5cm}|p{3cm}|p{2.8cm}|p{2.5cm}|p{2.7cm}|}
    \hline 
    \textbf{Autores} & \textbf{Título} & \textbf{Descripción} & \textbf{Ventajas} & \textbf{Desventajas} & \textbf{Diferencias con propuesta} \\
    \hline\hline
    
    Ahmed Elmagdy, Walid Magdy, Tarek El-Ganainy & Fake It Till You Make It: Profiling Fake Users on PornHub & Este estudio se enfoca en detectar perfiles falsos dentro de la plataforma PornHub. Utiliza aprendizaje automático sobre datos de perfil, comentarios y metainformación. & Se basa en datos reales y masivos. Proporciona un marco de análisis para detección de identidad falsa. & Restringido a una sola plataforma. No aborda mensajes en tiempo real. & • Análisis en tiempo real vs. atributos estáticos \newline • Procesamiento local en Android \newline • Enfoque en mensajería vs. perfiles \\
    \hline
    
    Rui Cao, Roy Ka-Wei Lee, Tuan-Anh Hoang & DeepHate: Hate Speech Detection via Multi-Faceted Text Representations & DeepHate combina múltiples representaciones de texto para detectar discurso de odio en redes sociales usando redes neuronales. & Utiliza múltiples representaciones textuales. Supera métodos clásicos en clasificación de odio. & Requiere recursos computacionales significativos. No aborda escenarios móviles. & • Modelos livianos vs. modelos pesados \newline • Procesamiento local vs. servidor \newline • Catfishing/extorsión vs. discurso de odio \\
    \hline
    
    Matthew Edwards & Scam Survivors Sextortion Reports & Conjunto de datos con más de 41,000 testimonios de víctimas de sextorsión, incluyendo detalles sobre ofensores y métodos utilizados. & Proporciona datos reales sobre casos de sextorsión. Permite explorar patrones de comportamiento. & Puede contener sesgos de autoreporte. No diseñado para detección en tiempo real. & • Herramienta operativa vs. base de datos \newline • Detección inmediata vs. análisis retrospectivo \newline • Prevención automatizada vs. documentación \\
    \hline
    
    \end{tabularx}
\end{table}

\section{\label{ref-008}Propuesta de soluci{\'{o}}n%Edgar Diego Pina VargasEP2025-05-09T22:45:00Z0 Edgar Diego Pina VargasEP Edgar Diego Pina VargasEP Edgar Diego Pina Vargas Edgar Diego Pina Vargas EP EP 2025-05-09T22:45:00Z 2025-05-09T22:45:00Z 0 0 0 Es parte de la seccopm del  desarrollo, aquí se debe hablar sobre el propósito de cada uno de los elementos de la arquitectura, diferenciar el almacenamiento remoto del local, el propósito de cuentas por lado del servidor y el almacenamiento de mensajers o todo lo que sea requerido para que el llm trabaje de forma local totalmente off line Es parte de la seccopm del  desarrollo, aquí se debe hablar sobre el propósito de cada uno de los elementos de la arquitectura, diferenciar el almacenamiento remoto del local, el propósito de cuentas por lado del servidor y el almacenamiento de mensajers o todo lo que sea requerido para que el llm trabaje de forma local totalmente off line   Justificar el uso del servidor local para servir como nodo central del sistema de mensajería: sobre el requerimiento de sincronicidad entre clientes esta parte va en propuesta de soluciion   Justificar el uso del servidor local para servir como nodo central del sistema de mensajería: sobre el requerimiento de sincronicidad entre clientes esta parte va en propuesta de soluciion   Justificar el uso del servidor local para servir como nodo central del sistema de mensajería: sobre el requerimiento de sincronicidad entre clientes esta parte va en propuesta de soluciion        Justificar el uso del servidor local para servir como nodo central del sistema de mensajería: sobre el requerimiento de sincronicidad entre clientes esta parte va en propuesta de soluciion
}

\begin{figure}[H]
    \centering
    \includegraphics[width=0.5\linewidth]{imagenes/esquema1.png}
    \caption{Enter Caption}
    \label{fig:enter-label}
\end{figure}

\begin{figure}[H]
    \centering
    \includegraphics[width=0.5\linewidth]{imagenes/esquema2.png}
    \caption{Enter Caption}
    \label{fig:enter-label}
\end{figure}


Anteriormente un objetivo demasiado amplio -{\textgreater} mediante la ejecuci{\'{o}}n de pruebas de identificaci{\'{o}}n de catphising a trav{\'{e}}s de ejemplos adaptados al espa{\~{n}}ol y a jergas propias de ciberfraude. generados a partir de ChatGPT almacenados en documentos CSV a incluir dentro de prompts automatizados para definir el modelo inicial con el cual trabajar el finetuning posteriormente 


\foreignlanguage{spanish}{\textbf{\textit{Arquitectura general del sistema (explicaci{\'{o}}n)}}}

\foreignlanguage{spanish}{Implementaci{\'{o}}n con enfoque edge processing.}

\foreignlanguage{spanish}{La soluci{\'{o}}n involucra un enfoque que distribuye los procesos en los dispositivos m{\'{o}}viles, para ello se contempla el desarrollo de una aplicaci{\'{o}}n nativa en Android sobre la implementaci{\'{o}}n en aplicaciones web progresivas PWA, por el nivel de integraci{\'{o}}n con el hardware}

\foreignlanguage{spanish}{Dentro de la aplicaci{\'{o}}n se gestionar{\'{i}}an diferentes servicios o m{\'{o}}dulos como lo son:}
\begin{itemize}
\item \foreignlanguage{spanish}{La ejecuci{\'{o}}n del LLM}
\item \foreignlanguage{spanish}{Gestor de notificaciones}
\item \foreignlanguage{spanish}{Gestor de bases de datos}
\item \foreignlanguage{spanish}{Api de conexi{\'{o}}n con el servidor centralizado}
\end{itemize}
\foreignlanguage{spanish}{En el desarrollo se contempla la implementaci{\'{o}}n de 2 bases de datos una local y otra centralizada, es importante diferenciar los prop{\'{o}}sitos de cada una, hablando de la base local esta recolectar{\'{i}}a informaci{\'{o}}n sobre las interacciones del usuario de manera estrictamente local en conformidad con la Ley Federal de Protecci{\'{o}}n de Datos Personales en Posesi{\'{o}}n de los Particulares (LFPDPPP) la privacidad del usuario}

\foreignlanguage{spanish}{\textbf{\textit{Tecnolog{\'{i}}as contempladas en la propuesta}}}

\foreignlanguage{spanish}{\textbf{Pruebas de desempe{\~{n}}o de identificaci{\'{o}}n en llms base }}\foreignlanguage{spanish}{(}\foreignlanguage{spanish}{\textbf{entorno de pruebas}}\foreignlanguage{spanish}{): se contempla el uso de t{\'{e}}cnicas de }\foreignlanguage{spanish}{\textit{prompt engineering}}\foreignlanguage{spanish}{ para maximizar la capacidad de modelos de lenguaje (LLM) en la tarea de clasificaci{\'{o}}n de mensajes sospechosos.}

\foreignlanguage{spanish}{Ollama: para ejecutar modelos localmente en el entorno de pruebas}

\foreignlanguage{spanish}{Jypiter: de Anaconda para automatizar la comunicaci{\'{o}}n con el modelo desde Python hacia Ollama}

\foreignlanguage{spanish}{\textbf{Implementaci{\'{o}}n del lado del servidor}}

\foreignlanguage{spanish}{PostgreSQL: sistema de gesti{\'{o}}n de bases de datos relacional utilizado para almacenar registros de conversaci{\'{o}}n procesados, configuraciones personalizadas por usuario y resultados de inferencia.}

\foreignlanguage{spanish}{Django: framework web de alto nivel en Python, empleado para construir las APIs que permiten la comunicaci{\'{o}}n entre la aplicaci{\'{o}}n de mensajer{\'{i}}a y el servidor.}

\foreignlanguage{spanish}{Azure o Ngrok: seg{{\'{u}}}n el entorno de despliegue, el servidor puede ser alojado bajo un esquema on-premise mediante t{{\'{u}}}neles seguros con Ngrok, o bien, mediante una infraestructura como servicio (IaaS) en la nube a trav{\'{e}}s de Microsoft Azure, lo que brinda escalabilidad y disponibilidad remota en entornos reales.}

\textbf{Aplicaci{\'{o}}n de Android}: \foreignlanguage{spanish}{El sistema de detecci{\'{o}}n se integra en una aplicaci{\'{o}}n m{\'{o}}vil desarrollada con Android, permitiendo an{\'{a}}lisis de mensajes directamente en el dispositivo del usuario. }

\foreignlanguage{spanish}{LiteRT (TensorFlow Lite): tecnolog{\'{i}}a que permite ejecutar modelos de lenguaje previamente ajustados en dispositivos Android, sin necesidad de conexi{\'{o}}n a internet o procesamiento en la nube. Esta herramienta es clave para el cumplimiento de principios de privacidad y para brindar respuestas en tiempo real.}

\foreignlanguage{spanish}{Android Studio: entorno de desarrollo integrado (IDE) utilizado para construir la interfaz de usuario, l{\'{o}}gica de comunicaci{\'{o}}n con el modelo y el servidor, y todas las funcionalidades necesarias para una aplicaci{\'{o}}n de mensajer{\'{i}}a segura.}

\section{Desarrollo}

Pruebas iniciales a modelos largos de lenguaje locales 

Para la generación de datos sintéticos se opta por un enfoque few shot para la formulación del prompt que se ocupara para
dar la instrucción a un modelo como Chat GPT o deepseek con madurez en tareas de generación de texto, se observo que
la generación sin datos puede dar lugar a datos poco realistas y poco variables por lo que en afinidad con las 8 practicas de
generación sintetica que plantea IBM en la sección del marco teorico se opta por few shot por sobre zero shot

https://huggingface.co/datasets/TExtPhish/TExtPhish
https://www.csusm.edu/iits/services/security/security-guidance/phishing_extortion.html
https://safecomputing.umich.edu/protect-yourself/phishing-scams/common-scams/sextortion
https://www.ic3.gov/PSA/2016/PSA160601

\section{Pruebas y resultados}

\section{\label{ref-010}Conclusiones}

\section{\label{ref-011}Trabajo a futuro%Edgar Diego Pina VargasEP2025-05-08T20:01:00Z0 Edgar Diego Pina VargasEP Edgar Diego Pina VargasEP Edgar Diego Pina Vargas Edgar Diego Pina Vargas EP EP 2025-05-08T20:01:00Z 2025-05-08T20:01:00Z 0 0 0   Trabajo a futuro podría implicar trabajar colaborativamente con estudiantes de leyes   Trabajo a futuro podría implicar trabajar colaborativamente con estudiantes de leyes   Trabajo a futuro podría implicar trabajar colaborativamente con estudiantes de leyes        Trabajo a futuro podría implicar trabajar colaborativamente con estudiantes de leyes
}

\section{\label{ref-012}Referencias}

\subsubsection{Introducci{\'{o}}n}

Usuarios de internet, por entidad federativa, seg{{\'{u}}}n principales usos, 2023: \url{https://view.officeapps.live.com/op/view.aspx?src=https\%3A\%2F\%2Fwww.inegi.org.mx\%2Fcontenidos\%2Fprogramas\%2Fendutih\%2F2023\%2Ftabulados\%2F2023_u2ed244.xlsx&wdOrigin=BROWSELINK}

Tasa de incidencia delictiva por cada cien mil habitantes por tipo de delito.

\url{https://www.inegi.org.mx/app/tabulados/interactivos/?px=ENVIPE_INCIDENCIA_TIPO_DELITO&bd=ENVIPE}

\subsection{Glosario de t{\'{e}}rminos}

IBM. (s.\,f.). \textit{{\textquestiondown}Qu{\'{e}} es la ingenier{\'{i}}a de prompts?} IBM Think. \url{https://www.ibm.com/es-es/think/topics/prompt-engineering}

IBM. (s.\,f.). \textit{{\textquestiondown}Qu{\'{e}} es el fine-tuning?} IBM Think. \url{https://www.ibm.com/es-es/think/topics/fine-tuning}

\subsection{Estado del arte}

Davidson, T., Warmsley, D., Macy, M., \& Weber, I. (2017). \textit{Automated Hate Speech Detection and the Problem of Offensive Language}. arXiv. \url{https://arxiv.org/abs/1705.06530}

Gupta, M., Zafar, H., \& Shukla, A. (2025). \textit{Multilingual Hate Speech Detection in Code-Mixed Data Using Prompt-Tuning}. arXiv. \url{https://arxiv.org/abs/2501.05443}

Alam, F., Shaar, S., Mubarak, H., \& Nakov, P. (2022). \textit{Large-scale multilingual hate speech detection dataset from social media} {[}Data set{]}. Mendeley Data. \url{https://doi.org/10.17632/xp8n69sdtp.1}

\subsection{Marco teorico}

Amazon Web Services. (s.f.). \textit{What is a large language model (LLM)?} \url{https://aws.amazon.com/what-is/large-language-model/}

Google for Developers. (s.f.). \textit{LLM: Ajuste, destilaci{\'{o}}n e ingenier{\'{i}}a de instrucciones}. https://developers.google.com/machine-learning/llm-foundations/llm

He, J., Wang, C., Liu, J., Xie, Y., Liu, P., Zheng, C., ... \& Liu, Z. (2023). \textit{A survey on large multimodal models}. arXiv. \url{https://arxiv.org/abs/2312.01044}

IBM. (s.f.). \textit{Large language models}. \url{https://www.ibm.com/think/topics/large-language-models}

Vaswani, A., Shazeer, N., Parmar, N., Uszkoreit, J., Jones, L., Gomez, A. N., ... \& Polosukhin, I. (2017). \textit{Attention is all you need}. arXiv. \url{https://arxiv.org/abs/1706.03762}

\subsubsection{Definici{\'{o}}n formal de lenguaje}

Tellez, E. S., Moctezuma, D., Miranda, S., Graff, M., \& Ruiz, G. (2022). \textit{Regionalized models for Spanish language variations based on Twitter}. arXiv. \url{https://doi.org/10.48550/arXiv.2110.06128}

\foreignlanguage{english}{Gonz{\'{a}}lez, J. A., \& Fern{\'{a}}ndez, M. (2018). Formal approaches to language: A review. }\foreignlanguage{english}{\textit{Journal of Applied Linguistics}}\foreignlanguage{english}{, 45(3), 123\textendash{}135. }\url{https://doi.org/10.1016/j.jal.2018.07.005}

\subsubsection{\foreignlanguage{spanish}{Propuesta de soluci{\'{o}}n}}

Google. (s.f.). TensorFlow Lite para entornos Edge (LiteRT). Google AI. \url{https://ai.google.dev/edge/litert?hl=es-419}

PostgreSQL Global Development Group. (s.f.). PostgreSQL: The World's Most Advanced Open Source Relational Database. https://www.postgresql.org/Ollama, Inc. (s.f.). 

Ollama: Run large language models locally. \url{https://ollama.com/}

Anaconda, Inc. (s.f.). Jupyter \textendash{} Anaconda Documentation. https://anaconda.org/anaconda/jupyter

\subsection{Desarrollo}

https://www.ibm.com/think/insights/synthetic-data-generation
\end{document}
